\section{Problem Definition}
\label{sec:problem}

\subsection{Rating Prediction Task}

The core problem addressed by this system is \textit{rating prediction}: given a user $u$ and an item 
(movie) $i$, predict the numerical rating $r_{u,i}$ that the user would assign to the item. Formally:

\begin{equation}
  f: (u, i) \mapsto \hat{r}_{u,i} \in [1, 5]
\end{equation}

where:
\begin{itemize}
  \item $u \in \mathcal{U}$ is a user identifier from the set of all users.
  \item $i \in \mathcal{I}$ is an item (movie) identifier from the set of all items.
  \item $\hat{r}_{u,i}$ is the predicted rating, constrained to the interval $[1.0, 5.0]$.
\end{itemize}

The prediction is made using historical rating data $\mathcal{R} = \{(u_j, i_j, r_{u_j, i_j})_{j=1}^{|\mathcal{R}|}\}$, 
representing past observations of users' ratings on movies.

\subsection{Setting and Assumptions}

We operate under the following setting:

\textbf{Implicit Feedback Context:} While the MovieLens 100K dataset contains explicit numerical ratings, 
the system must handle both known and unknown user-item pairs. For unknown pairs (users or movies not in the 
training set), the model is expected to produce reasonable estimates based on learned latent factors.

\textbf{Cold Start Simplifications:} We do not implement sophisticated cold-start strategies. Instead, we 
rely on the base model's ability to extrapolate from latent factor space. This is a practical limitation 
addressed in Section~\ref{sec:limitations}.

\textbf{Static Model Assumption:} The model is trained once and served in production without online learning. 
This is typical for batch-trained systems and simplifies operational deployed.

\textbf{Single Prediction Paradigm:} Beyond single-prediction requests, the API also supports batch 
predictions---processing multiple $(u, i)$ pairs in a single request. Both modes share the same underlying 
prediction logic.

\subsection{Formal Data Representation}

The rating data is typically represented as a sparse matrix $R \in \mathbb{R}^{|\mathcal{U}| \times |\mathcal{I}|}$, 
where entry $R_{u,i}$ is the rating given by user $u$ to movie $i$, or undefined if the rating was never 
provided. This sparsity is fundamental to the problem: real-world user-item interactions are sparse, 
and any effective recommender system must handle this sparsity gracefully.

\subsection{Success Criteria}

The system is evaluated against the following criteria:

\begin{enumerate}
  \item \textbf{Correctness}: Predictions are real-valued and fall within $[1.0, 5.0]$.
  
  \item \textbf{Reliability}: The system handles edge cases (unknown users, missing data) without crashing.
  
  \item \textbf{Performance}: Prediction latency is acceptable (sub-second) for serving via REST API.
  
  \item \textbf{Usability}: The API is well-documented, validates input, and returns meaningful error messages.
  
  \item \textbf{Operability}: The system can be deployed, monitored, and scaled using standard DevOps practices.
\end{enumerate}

These criteria span model accuracy, system robustness, and operational concerns---reflecting the reality 
that production ML systems are judged on far more than raw prediction accuracy.
